% ============================================================
% Extended related work (appendix). Lines that didn't make the
% main paper §\ref{sec:related} because they are either:
%   (a) easy-to-confuse but methodologically distant (RAG/corpus),
%   (b) at a different layer (LLM text watermarking), or
%   (c) standard cryptographic plumbing (commitments, Merkle trees).
% ============================================================
\section{Extended Related Work}
\label{app:related-ext}

% ============================================================
\subsection{RAG / Corpus-Level Watermarking}
\label{app:related-rag}
% ============================================================
\todo{Cite RAG-WM~\cite{ragwm}, Ward / RAG-DI~\cite{ragdi},
Watermarked Canaries~\cite{canaries}, AQUA~\cite{aqua},
\kgmark{}~\cite{kgmark}, Graph DB Pseudo-Nodes~\cite{graphdbwm}.
Common pattern: protect dataset / corpus IP by injecting watermarked
content into a static knowledge base, then detect via black-box LLM
query.}

\todo{Difference to \memmark{}: (i) target — those protect data IP
("who stole my documents"); \memmark{} does writer attribution
("which agent evolved this memory state"). (ii) mechanism — static
content injection vs behavioral keyed selection on evolve decisions.
(iii) verification — those need LLM black-box query at detection
time; \memmark{} verifies in-record from snapshot + key
(\S\ref{sec:rq3} R3).}

\todo{Among them, \kgmark{} is the closest single-point baseline and
appears in the main paper's experiments — see
\S\ref{sec:related-memory}. The rest of this RAG/corpus line is
mechanistically distinct enough that we don't carry it through
\S\ref{sec:exp}.}

% ============================================================
\subsection{LLM Text Watermarking}
\label{app:related-llmtext}
% ============================================================
\todo{Cite KGW~\cite{kgw}, SynthID-Text~\cite{synthid},
MarkLLM~\cite{markllm}. Token-level watermarks on the final LLM
output text. Different layer entirely: \memmark{} does not touch
output text, only the memory backend's internal evolve decisions.}

\todo{Composability — \memmark{} is orthogonal to text-layer
watermarks. \S\ref{sec:rq6} / Appendix~\ref{app:composability}
verifies that running both simultaneously preserves each one's
detection ability without joint utility collapse.}

% ============================================================
\subsection{Cryptographic Provenance Primitives}
\label{app:related-crypto}
% ============================================================
\todo{Commit-then-reveal commitment schemes, Merkle trees and
transparency logs (e.g., Certificate Transparency~\cite{ct}),
PRF-keyed nonces. All are standard cryptographic primitives.}

\todo{Contribution at this layer: \memmark{}'s \S\ref{sec:method-audit}
composes them into a memory-evolve-aware audit trace
(per-decision commitment + per-session Merkle log + signed
in-snapshot anchor + per-leaf inclusion proof). The novelty is
systems-level: how these primitives are wired into a third-party
memory SDK without modifying it, and how they collectively support
the R1 / R2 / R3 verification hierarchy.}
